\section{附录}
\subsection{不同日期的轨迹结果}
\label{analysis:reuse-evolution}

\begin{figure}[H]
    \centering
    \includegraphics[width=\columnwidth]{./data/reuse-time-probability-distribution-7days-2\figsuffix}
    \caption{5 个工作日中复用概率累积分布函数的演化。第一行：上午 9:00 至 10:00 的分布；第二行：上午 10:00 至 11:00 的分布。}
    \label{fig:motiv-7days}
\end{figure}

\fig{fig:motiv-7days}展示了连续 5 个工作日中，各请求类别在未来时间窗口内的复用概率累积分布函数（CDF）如何演化。可以观察到，工作日上午 9:00 至 10:00 的复用概率彼此相似。因此，可以离线分析前一天的复用日志，预测次日各类别的复用概率。不过，图像类别的复用概率比其他类别更难预测；作者推测，这是因为图像类别的可复用粒度比文本类别更粗，而且不同用户之间的差异更大。
